\section{Conceptual Model: Music as a Recursive Hierarchy}\label{sec:conceptual-model}

The fundamental premise of this research is that musical composition is inherently structured and hierarchical.
In traditional music theory, a piece is not merely a sequence of notes, but a complex organization of motifs, phrases, sections, and movements.
However, many modern digital tools and domain-specific languages treat music as a linear stream of events or a flat sequence of commands.
This research proposes a shift toward a \textit{structural representation}, where musical ideas are treated as modular, reusable, and composable units.

\subsection{Linear vs. Structural Representation}

In linear representation models, commonly found in live coding environments, the focus is on the \textit{performance}.
Time is often modeled as a continuous thread where commands such as \texttt{play} and \texttt{sleep} are executed sequentially to generate sound.
While effective for improvisation and real-time interaction, this model often sacrifices the ability to express complex temporal relationships and hierarchical dependencies.
Modifying a single motif in a linear script often requires broad changes to subsequent timing logic, leading to fragile and non-modular code.

In contrast, the proposed structural representation models music as a \textit{recursive hierarchy}.
Under this model, a composition is built by nesting smaller musical units within larger ones.
A single note is the base case of this hierarchy.
Notes are grouped into motifs (patterns), motifs are sequenced into phrases, and phrases are assigned to parallel execution contexts (tracks).
This approach allows for a clean separation of concerns: the musical content (the "what") is defined independently of its temporal placement (the "when").

\subsection{Music as a Pure Function}

By treating musical structures as first-class entities, the language allows them to be manipulated using concepts derived from functional programming.
A musical pattern can be viewed as a \textit{pure function}: it encapsulates a specific musical logic that produces a deterministic sequence of events.
Just as a function in programming can be called multiple times with different arguments, a musical pattern in the proposed DSL can be instantiated across different tracks, at different times, or with different parameters (such as pitch transpositions or rhythmic stretching).

This functional mapping provides several advantages:
\begin{itemize}
    \item \textbf{Modularity:} Musical ideas can be defined once and reused throughout a composition.
    \item \textbf{Abstraction:} Complex polyphonic structures can be hidden behind simple pattern interfaces.
    \item \textbf{Determinism:} Because the language lacks a real-time runtime, every musical relationship is resolved at compile-time, ensuring a "frozen" and perfect mathematical representation in the output artifact.
\end{itemize}
